\section{The CPU Execution Cycle}
Once the processes are loaded into the ready queue, the CPU threads execute the \texttt{cpu\_routine}. This routine acts as a perpetual Fetch-Decode-Execute cycle.

\subsection{Process Dispatching}
The CPU constantly polls the scheduler for work.
\begin{lstlisting}[language=C, caption=Process Dispatch in \texttt{cpu\_routine} (\texttt{os.c})]
while (1) {
    if (proc == NULL) {
        proc = get_proc(); // Fetches highest priority process
        if (proc == NULL) {
            next_slot(timer_id);
            continue; 
        }
    } else if (proc->pc == proc->code->size) {
        // Process finished executing all instructions
        printf("\tCPU %d: Processed %2d has finished\n", id, proc->pid);
        remove_running_proc(proc);
        free(proc);
        proc = get_proc();
        time_left = 0;
    } else if (time_left == 0) {
        // Time slice exhausted - Preemption
        printf("\tCPU %d: Put process %2d to run queue\n", id, proc->pid);
        put_proc(proc); // Put back to its priority queue
        proc = get_proc();
    }
    
    // Recheck: if no process available and loader is done, stop CPU
    if (proc == NULL && done) {
        printf("\tCPU %d stopped\n", id);
        break;
    } else if (proc == NULL) {
        next_slot(timer_id);
        continue;
    } else if (time_left == 0) {
        printf("\tCPU %d: Dispatched process %2d\n", id, proc->pid);
        time_left = time_slot;
    }
    
    // Execute one instruction
    run(proc);
    time_left--;
    next_slot(timer_id);
}
\end{lstlisting}
This loop implements \textbf{Preemptive Scheduling}. A process cannot hold the CPU forever because \texttt{time\_left} is decremented every cycle. Once it hits zero, the CPU forces a context switch by putting the process back in the queue and fetching the next one.

\begin{theorybox}[Process Termination Path]
When a process finishes all its instructions (\texttt{proc->pc == proc->code->size}), the CPU must safely remove it from the system. Rather than calling \texttt{purgequeue()} directly on the \texttt{running\_list}---which would be unsafe in a multi-threaded environment---the code calls the wrapper function \texttt{remove\_running\_proc(proc)} (defined in \texttt{sched.c}). This wrapper acquires the global \texttt{queue\_lock} mutex before invoking \texttt{purgequeue(\&running\_list, proc)}, ensuring that concurrent CPU threads do not corrupt the list during removal:
\begin{lstlisting}[language=C, caption=Thread-safe process removal (\texttt{sched.c})]
void remove_running_proc(struct pcb_t *proc) {
    pthread_mutex_lock(&queue_lock);
    purgequeue(&running_list, proc);
    pthread_mutex_unlock(&queue_lock);
}
\end{lstlisting}
\end{theorybox}

\subsection{Instruction Execution (\texttt{run()})}
The \texttt{run()} function decodes the instruction pointed to by the Program Counter (\texttt{pc}) and routes it to the corresponding OS subsystem. It returns an \texttt{int} status code: 0 indicates successful execution and 1 signals a failure (e.g., invalid register index, denied kernel-mode instruction).

\begin{lstlisting}[language=C, caption=Instruction Routing in \texttt{cpu.c}]
int run(struct pcb_t *proc) {
    if (proc->pc >= proc->code->size) return 1;

    struct inst_t ins = proc->code->text[proc->pc];
    proc->pc++;
    int stat = 1;
    switch (ins.opcode) {
    case CALC:
        stat = calc(proc);
        break;
    case ALLOC:
#ifdef MM_PAGING
        stat = liballoc(proc, ins.arg_0, ins.arg_1);
#else
        stat = alloc(proc, ins.arg_0, ins.arg_1);
#endif
        break;
    case KMALLOC:
        if (proc->mode_bit == 1) { stat = 1; break; }
        stat = libkmem_malloc(proc, ins.arg_0, ins.arg_1);
        break;
    case FREE:
#ifdef MM_PAGING
        stat = libfree(proc, ins.arg_0);
#else
        stat = free_data(proc, ins.arg_0);
#endif
        break;
    case READ:
#ifdef MM_PAGING
        {
            uint32_t read_val = 0;
            stat = libread(proc, ins.arg_0, ins.arg_1, &read_val);
            if (stat == 0) {
                if (valid_reg_index(proc, ins.arg_2))
                    proc->regs[ins.arg_2] = read_val;
                else stat = 1;
            }
        }
#else
        stat = read(proc, ins.arg_0, ins.arg_1, ins.arg_2);
#endif
        break;
    case WRITE:
#ifdef MM_PAGING
        stat = libwrite(proc, ins.arg_0, ins.arg_1, ins.arg_2);
#else
        stat = write(proc, ins.arg_0, ins.arg_1, ins.arg_2);
#endif
        break;
    case SYSCALL:
        proc->mode_bit = 0;  // Enter Kernel Mode
        stat = libsyscall(proc, ins.arg_0, ins.arg_1, ins.arg_2, ins.arg_3);
        proc->mode_bit = 1;  // Return to User Mode
        break;
    default:
        stat = 1;
    }
    return stat;
}
\end{lstlisting}

\begin{infobox}[Dual-Mode Operation \& Mode Bit Transitions]
The \texttt{SYSCALL} case in the switch statement illustrates the simulator's dual-mode protection mechanism. Before dispatching the system call, the CPU sets \texttt{proc->mode\_bit = 0} (Kernel Mode), and after the system call returns, it restores \texttt{proc->mode\_bit = 1} (User Mode). This sandwich ensures that \texttt{libsyscall} and the kernel handler it invokes execute with full kernel privileges, while all other instructions run in restricted user mode.

Crucially, kernel-only instructions (\texttt{KMALLOC}, \texttt{KMEM\_CACHE\_CREATE}, \texttt{KMEM\_CACHE\_ALLOC}, \texttt{COPY\_FROM\_USER}, \texttt{COPY\_TO\_USER}) are \textbf{guarded} by an explicit \texttt{if (proc->mode\_bit == 1) \{ stat = 1; break; \}} check. Since processes always start with \texttt{mode\_bit = 1} (set in \texttt{loader.c:59}), these instructions are silently rejected unless the process is currently inside a system call handler. This mirrors the hardware protection ring architecture of real CPUs.
\end{infobox}

\begin{infobox}[Virtual Registers \& Symbol Table]
The simulated architecture provides \textbf{10 general-purpose registers} per process (\texttt{regs[0]}--\texttt{regs[9]}), as defined in \texttt{common.h}:
\begin{lstlisting}[language=C]
addr_t regs[10];  // Registers, store address of allocated regions
\end{lstlisting}
The validation function \texttt{valid\_reg\_index()} ensures instructions don't overflow this array. \texttt{CALC} operates entirely in registers. Memory operations like \texttt{ALLOC} return virtual addresses stored in these registers for later use by \texttt{READ}/\texttt{WRITE}.

Note: The \texttt{mm\_struct} separately maintains a \textbf{symbol table} (\texttt{symrgtbl[30]}) with 30 entries (\texttt{PAGING\_MAX\_SYMTBL\_SZ = 30}) that maps register indices to full virtual memory region descriptors (\texttt{vm\_rg\_struct}). This is a kernel-side tracking structure, distinct from the 10 CPU registers.
\end{infobox}